\documentclass[11pt, a4paper]{article}

% bfw-style.tex — gemeinsame Style-Definitionen für alle Handouts/Aufgabenhefte
% Voraussetzung: Kompilation mit xelatex (wegen fontspec)

\usepackage[ngerman]{babel}
\usepackage{geometry}
\geometry{a4paper, top=2.2cm, bottom=2.2cm, left=2.3cm, right=2.3cm}

\usepackage{fontspec}
% Fallback-Kette: Source Sans Pro -> Calibri -> Segoe UI -> Latin Modern Sans
% (Latin Modern ist immer mit MiKTeX vorhanden)
\IfFontExistsTF{Source Sans Pro}{%
  \setmainfont{Source Sans Pro}[Ligatures=TeX]%
}{\IfFontExistsTF{Calibri}{%
  \setmainfont{Calibri}[Ligatures=TeX]%
}{\IfFontExistsTF{Segoe UI}{%
  \setmainfont{Segoe UI}[Ligatures=TeX]%
}{%
  \setmainfont{Latin Modern Sans}[Ligatures=TeX]%
}}}
\IfFontExistsTF{Source Code Pro}{%
  \setmonofont{Source Code Pro}[Scale=0.92]%
}{\IfFontExistsTF{Consolas}{%
  \setmonofont{Consolas}[Scale=0.92]%
}{%
  \setmonofont{Latin Modern Mono}[Scale=0.92]%
}}

\usepackage{xcolor}
% Farbpalette: hell, freundlich, modern
\definecolor{bfwPrimary}{HTML}{0EA5A4}    % Petrol/Türkis - Primär
\definecolor{bfwPrimaryDark}{HTML}{0F766E} % dunkler Petrol für Headings
\definecolor{bfwAccent}{HTML}{F59E0B}     % Amber - Akzent
\definecolor{bfwSuccess}{HTML}{10B981}    % Grün - Erfolg / Lösung
\definecolor{bfwWarn}{HTML}{F97316}       % Orange - Achtung
\definecolor{bfwInk}{HTML}{0F172A}        % fast Schwarz
\definecolor{bfwInkSoft}{HTML}{475569}    % gedämpftes Grau
\definecolor{bfwBgSoft}{HTML}{F8FAFC}     % off-white
\definecolor{bfwBgTeal}{HTML}{ECFEFF}     % helles Türkis
\definecolor{bfwBgAmber}{HTML}{FEF3C7}    % helles Gelb
\definecolor{bfwBgRose}{HTML}{FFE4E6}     % helles Rosé für Eselsbrücken

\usepackage[colorlinks=true, linkcolor=bfwPrimaryDark, urlcolor=bfwPrimaryDark]{hyperref}
\usepackage{enumitem}
\setlist[itemize]{leftmargin=*, itemsep=2pt, topsep=2pt}
\setlist[enumerate]{leftmargin=*, itemsep=2pt, topsep=2pt}

\usepackage[most]{tcolorbox}
\usepackage{booktabs}
\usepackage{tikz}
\usetikzlibrary{decorations.pathmorphing, positioning, shapes.geometric, arrows.meta}

% Merkkästchen — wichtige Definition / Kernpunkt
\newtcolorbox{merksatz}[1][]{%
  enhanced, breakable,
  colback=bfwBgTeal,
  colframe=bfwPrimary,
  boxrule=0.6pt,
  arc=4pt,
  left=10pt, right=10pt, top=8pt, bottom=8pt,
  fonttitle=\bfseries\color{bfwPrimaryDark},
  title={\faLightbulb\ Merksatz},
  #1
}

% Eselsbrücke — Mnemoniks
\newtcolorbox{eselsbruecke}[1][]{%
  enhanced, breakable,
  colback=bfwBgRose,
  colframe=bfwAccent,
  boxrule=0.6pt,
  arc=4pt,
  left=10pt, right=10pt, top=8pt, bottom=8pt,
  fonttitle=\bfseries\color{bfwWarn},
  title={Eselsbrücke},
  #1
}

% Beispiel-Box
\newtcolorbox{beispiel}[1][]{%
  enhanced, breakable,
  colback=bfwBgSoft,
  colframe=bfwInkSoft,
  boxrule=0.4pt,
  arc=3pt,
  left=10pt, right=10pt, top=8pt, bottom=8pt,
  fonttitle=\bfseries\color{bfwInk},
  title={Beispiel},
  #1
}

% Lösungs-Box (für Aufgabenhefte)
\newtcolorbox{loesung}[1][]{%
  enhanced, breakable,
  colback=bfwBgAmber!40,
  colframe=bfwSuccess,
  boxrule=0.6pt,
  arc=3pt,
  left=10pt, right=10pt, top=8pt, bottom=8pt,
  fonttitle=\bfseries\color{bfwSuccess!70!black},
  title={Lösung},
  #1
}

% Glossar-Eintrag
\newcommand{\glossareintrag}[2]{%
  \par\noindent\textbf{\color{bfwPrimaryDark}#1} \enspace #2\par\smallskip
}

% Section-Stil — etwas farbiger als Standard
\usepackage{titlesec}
\titleformat{\section}
  {\Large\bfseries\color{bfwPrimaryDark}}
  {\thesection}{0.6em}{}
\titleformat{\subsection}
  {\large\bfseries\color{bfwPrimary}}
  {\thesubsection}{0.5em}{}
\titleformat{\subsubsection}
  {\normalsize\bfseries\color{bfwInk}}
  {\thesubsubsection}{0.4em}{}

% TikZ-Stil "skizzenhaft" — etwas weicher als Rough.js, aber stabil
\tikzset{
  roughbox/.style = {
    draw=bfwPrimaryDark, thick, fill=bfwBgTeal,
    rounded corners=4pt, inner sep=6pt
  },
  roughaccent/.style = {
    draw=bfwAccent, thick, fill=bfwBgAmber,
    rounded corners=4pt, inner sep=6pt
  },
  roughline/.style = {
    draw=bfwInkSoft, thick, -{Stealth[length=2.5mm]}
  }
}

% Bullet-Symbole (bewusst kein FontAwesome — vermeidet Font-Installations-Probleme)
\newcommand{\faLightbulb}{$\bullet$}
\newcommand{\faBookmark}{$\bullet$}

% Header/Footer
\usepackage{fancyhdr}
\pagestyle{fancy}
\fancyhf{}
\renewcommand{\headrulewidth}{0pt}
\fancyfoot[L]{\small\color{bfwInkSoft}\thepage}
\fancyfoot[R]{\small\color{bfwInkSoft} BFW M-064 Beschaffung im Büromanagement}


\title{\vspace*{-1.5cm}\Huge\bfseries\color{bfwPrimaryDark}Aufgabenheft Tag~9\\[0.4em]
  \large\color{bfwInkSoft}\textnormal{Kommunikationsformen \& Teamarbeit --- Übungen im IHK-Klausurformat}}
\author{\textsf{Büroprozesse gestalten und Arbeitsvorgänge organisieren \textperiodcentered{} M-067}\\
\textsf{\small\copyright\ Can Siebert\,\textperiodcentered\,2026}}
\date{Selbstlernmaterial \textperiodcentered{} 09.07.2026}

\begin{document}
\maketitle
\thispagestyle{empty}

\begin{merksatz}[title={\bfseries So nutzen Sie dieses Aufgabenheft}]
Bearbeiten Sie alle Aufgaben zunächst \emph{ohne} in die Lösungen zu schauen. Schreiben
Sie Ihre Antworten in vollständigen Sätzen --- die IHK-Prüfung verlangt formulierte
Begründungen, nicht bloße Stichworte. Beachten Sie die \emph{Operatoren} (nennen,
erläutern, analysieren, beurteilen) --- sie geben den geforderten Antwort-Tiefegrad vor.
Erst danach öffnen Sie den Lösungsteil ab Seite~\pageref{loesungen} und vergleichen.
\end{merksatz}

\tableofcontents
\vspace{1em}
\hrule

\section{Aufgaben}

\subsection*{Aufgabe~1 --- Sender-Empfänger-Modell (10 Punkte)}

\noindent\textbf{\color{bfwAccent}YouTube-Vorbereitung:}\enspace \href{https://www.youtube.com/results?search_query=Sender+Empfaenger+Modell+einfach+erklaert}{Sender-Empfänger-Modell einfach erklärt} (\url{https://www.youtube.com/results?search_query=Sender+Empfaenger+Modell+einfach+erklaert})

\medskip
Im Großraumbüro der \emph{Lindner \& Partner GmbH} kommt es immer wieder zu
Missverständnissen zwischen den Beschäftigten.

\begin{enumerate}[label=(\alph*)]
  \item \textbf{Beschreiben} Sie den Ablauf einer Kommunikation nach dem Sender-Empfänger-Modell. \hfill (4~P.)
  \item \textbf{Erläutern} Sie den Begriff \emph{Rauschen} und nennen Sie zwei Beispiele aus dem Büro. \hfill (3~P.)
  \item \textbf{Beurteilen} Sie, warum Nachfragen keine Schwäche, sondern Qualitätssicherung ist. \hfill (3~P.)
\end{enumerate}

\subsection*{Aufgabe~2 --- Vier Seiten einer Nachricht (14 Punkte)}

\noindent\textbf{\color{bfwAccent}YouTube-Vorbereitung:}\enspace \href{https://www.youtube.com/results?search_query=Vier+Seiten+einer+Nachricht+Schulz+von+Thun}{Vier Seiten einer Nachricht} (\url{https://www.youtube.com/results?search_query=Vier+Seiten+einer+Nachricht+Schulz+von+Thun})

\medskip
Ein Vorgesetzter betritt das Büro und sagt zu einer Mitarbeiterin, die noch am Schreibtisch
sitzt: \emph{„Es ist schon nach 18 Uhr.”}

\begin{enumerate}[label=(\alph*)]
  \item \textbf{Nennen} Sie die vier Seiten einer Nachricht nach Schulz von Thun. \hfill (4~P.)
  \item \textbf{Analysieren} Sie die Aussage „Es ist schon nach 18 Uhr.” auf allen vier Ebenen. \hfill (6~P.)
  \item \textbf{Erläutern} Sie, wie ein Missverständnis entsteht, wenn die Mitarbeiterin mit dem \emph{Beziehungsohr} hört. \hfill (4~P.)
\end{enumerate}

\subsection*{Aufgabe~3 --- Verbal, nonverbal, paraverbal (12 Punkte)}

\noindent\textbf{\color{bfwAccent}YouTube-Vorbereitung:}\enspace \href{https://www.youtube.com/results?search_query=verbale+nonverbale+paraverbale+Kommunikation}{Verbale, nonverbale und paraverbale Kommunikation} (\url{https://www.youtube.com/results?search_query=verbale+nonverbale+paraverbale+Kommunikation})

\medskip
\begin{enumerate}[label=(\alph*)]
  \item \textbf{Erläutern} Sie den Unterschied zwischen verbaler, nonverbaler und paraverbaler Kommunikation. \hfill (3~P.)
  \item \textbf{Ordnen} Sie folgende Elemente den drei Ebenen \textbf{zu}: Wortwahl, Tonfall, Gestik, Lautstärke, Mimik, Sprechtempo. \hfill (3~P.)
  \item \textbf{Beurteilen} Sie das Problem, wenn verbale und nonverbale Botschaft nicht zusammenpassen. \hfill (3~P.)
  \item \textbf{Begründen} Sie, warum der Tonfall am Telefon besonders wichtig ist. \hfill (3~P.)
\end{enumerate}

\subsection*{Aufgabe~4 --- Kommunikationsformen \& E-Mail (15 Punkte)}

\noindent\textbf{\color{bfwAccent}YouTube-Vorbereitung:}\enspace \href{https://www.youtube.com/results?search_query=E-Mail+Knigge+Geschaeftsbrief+Regeln}{E-Mail-Knigge im Geschäftsverkehr} (\url{https://www.youtube.com/results?search_query=E-Mail+Knigge+Geschaeftsbrief+Regeln})

\medskip
Sie arbeiten in der Verwaltung und müssen für verschiedene Anlässe die passende
Kommunikationsform wählen.

\begin{enumerate}[label=(\alph*)]
  \item \textbf{Ordnen} Sie jedem Anlass die beste Kommunikationsform \textbf{zu} und \textbf{begründen} Sie kurz: \hfill (4~P.)
    \begin{enumerate}[label=\arabic*., leftmargin=*]
      \item vertrauliches Kritikgespräch mit einer Mitarbeiterin,
      \item verbindliche Bestellung mit Nachweis gegenüber einem Lieferanten,
      \item kurze Rückfrage an die Kollegin im Nebenraum,
      \item Projektbesprechung mit drei Standorten.
    \end{enumerate}
  \item \textbf{Nennen} Sie fünf Regeln des E-Mail-Knigge. \hfill (5~P.)
  \item \textbf{Erläutern} Sie, warum geschäftliche E-Mails Pflichtangaben enthalten müssen und wovon ihr Umfang abhängt. \hfill (4~P.)
  \item \textbf{Beurteilen} Sie, wann ein Messenger/Chat \emph{nicht} die richtige Wahl ist. \hfill (2~P.)
\end{enumerate}

\subsection*{Aufgabe~5 --- Teamarbeit (12 Punkte)}

\noindent\textbf{\color{bfwAccent}YouTube-Vorbereitung:}\enspace \href{https://www.youtube.com/results?search_query=Teamarbeit+Vorteile+Nachteile}{Teamarbeit --- Vorteile und Nachteile} (\url{https://www.youtube.com/results?search_query=Teamarbeit+Vorteile+Nachteile})

\medskip
\begin{enumerate}[label=(\alph*)]
  \item \textbf{Erläutern} Sie den Unterschied zwischen einer \emph{Gruppe} und einem \emph{Team}. \hfill (3~P.)
  \item \textbf{Nennen} Sie je drei Vor- und Nachteile der Teamarbeit. \hfill (6~P.)
  \item \textbf{Erläutern} Sie das \emph{soziale Faulenzen} und nennen Sie eine Gegenmaßnahme. \hfill (3~P.)
\end{enumerate}

\subsection*{Aufgabe~6 --- Teamphasen, Konflikt \& Feedback (15 Punkte)}

\noindent\textbf{\color{bfwAccent}YouTube-Vorbereitung:}\enspace \href{https://www.youtube.com/results?search_query=Teamphasen+Tuckman+Forming+Storming+Norming+Performing}{Teamphasen nach Tuckman} (\url{https://www.youtube.com/results?search_query=Teamphasen+Tuckman+Forming+Storming+Norming+Performing})

\medskip
Ein neues Projektteam aus fünf Personen wird zusammengestellt. In der zweiten Woche kommt
es zu Streit über die Aufgabenverteilung; einzelne Mitglieder ziehen sich zurück.

\begin{enumerate}[label=(\alph*)]
  \item \textbf{Nennen} Sie die fünf Teamphasen nach Tuckman in richtiger Reihenfolge und \textbf{beschreiben} Sie jede kurz. \hfill (5~P.)
  \item \textbf{Bestimmen} Sie, in welcher Phase sich das Team befindet, und \textbf{begründen} Sie. \hfill (4~P.)
  \item \textbf{Nennen} Sie vier Feedbackregeln, mit denen der Konflikt konstruktiv angesprochen werden kann. \hfill (4~P.)
  \item \textbf{Beurteilen} Sie die Aussage: \emph{„Streit im Team bedeutet, dass das Team gescheitert ist.”} \hfill (2~P.)
\end{enumerate}

\vspace{1em}
\noindent\textbf{Punkte gesamt: 78}

\newpage

\section{Lösungen}\label{loesungen}

\subsection*{Lösung Aufgabe~1}
\begin{loesung}
\textbf{(a)} Der \emph{Sender} hat eine Absicht und verschlüsselt sie in Zeichen
(\emph{Encodieren}: Worte, Tonfall, Körpersprache). Die \emph{Nachricht} wird über einen
\emph{Kanal} übertragen. Der \emph{Empfänger} nimmt sie auf und entschlüsselt sie
(\emph{Decodieren}) zu einer eigenen Bedeutung; anschließend gibt er ein \emph{Feedback}.

\textbf{(b)} \emph{Rauschen} ist jede Störung, die die Übertragung beeinträchtigt.
Beispiele: Lärm im Großraumbüro, schlechte Telefonverbindung, Ablenkung, unklare
Formulierung, unterschiedliches Vorwissen.

\textbf{(c)} Weil Encodieren und Decodieren subjektiv sind, sind Missverständnisse normal.
Nachfragen stellt sicher, dass die Nachricht richtig verstanden wurde --- es verhindert
Fehler und Doppelarbeit und ist damit aktive Qualitätssicherung, kein Zeichen von
Unwissen.
\end{loesung}

\subsection*{Lösung Aufgabe~2}
\begin{loesung}
\textbf{(a)} Die vier Seiten: \textbf{Sachebene, Selbstoffenbarung, Beziehungsebene,
Appellebene}.

\textbf{(b)} Analyse von „Es ist schon nach 18 Uhr.”:
\begin{itemize}[itemsep=0pt]
  \item \emph{Sache:} Die Uhrzeit ist nach 18 Uhr.
  \item \emph{Selbstoffenbarung:} „Ich achte auf die Arbeitszeit / mir fällt auf, dass du
    noch da bist.”
  \item \emph{Beziehung:} „Ich sorge mich um dich” --- oder „Du arbeitest unwirtschaftlich
    langsam.”
  \item \emph{Appell:} „Mach Feierabend” --- oder „Beeil dich.”
\end{itemize}

\textbf{(c)} Hört die Mitarbeiterin mit dem \emph{Beziehungsohr}, deutet sie die Aussage
als Kritik an ihrem Arbeitstempo („Du bist zu langsam”), obwohl der Vorgesetzte vielleicht
nur Fürsorge ausdrücken wollte. So entsteht aus einer sachlichen Aussage ein Konflikt.
\end{loesung}

\subsection*{Lösung Aufgabe~3}
\begin{loesung}
\textbf{(a)} \emph{Verbal} = der Inhalt, die Worte. \emph{Nonverbal} = Körpersprache ohne
Worte (Mimik, Gestik, Körperhaltung, Blickkontakt). \emph{Paraverbal} = die Art des
Sprechens (Tonfall, Lautstärke, Sprechtempo, Betonung, Pausen).

\textbf{(b)} Zuordnung:
\begin{itemize}[itemsep=0pt]
  \item \emph{Verbal:} Wortwahl.
  \item \emph{Nonverbal:} Gestik, Mimik.
  \item \emph{Paraverbal:} Tonfall, Lautstärke, Sprechtempo.
\end{itemize}

\textbf{(c)} Passen verbale und nonverbale Botschaft nicht zusammen (\emph{Inkongruenz}),
entsteht eine widersprüchliche Doppelbotschaft. Empfänger glauben im Zweifel der
Körpersprache und Stimme; die Glaubwürdigkeit leidet, es entstehen Verunsicherung und
Misstrauen.

\textbf{(d)} Am Telefon fehlt die nonverbale Ebene vollständig. Nur die Stimme
(\emph{paraverbal}) und die Worte übertragen die Botschaft --- Tonfall, Freundlichkeit und
Sprechtempo entscheiden deshalb besonders stark über den Eindruck beim Anrufer.
\end{loesung}

\subsection*{Lösung Aufgabe~4}
\begin{loesung}
\textbf{(a)} Zuordnung mit Begründung:
\begin{enumerate}[label=\arabic*., itemsep=0pt]
  \item Kritikgespräch: \textbf{persönliches Gespräch} --- vertraulich, Tonfall/Mimik
    wichtig, keine schriftliche Spur.
  \item Bestellung mit Nachweis: \textbf{E-Mail} (oder Brief) --- schriftlich,
    dokumentiert, beweissicher.
  \item kurze Rückfrage: \textbf{persönliches Gespräch oder Messenger/Chat} --- schnell und
    informell.
  \item Besprechung mit drei Standorten: \textbf{Videokonferenz} --- verbindet Standorte,
    spart Reisezeit, überträgt Bild und Ton.
\end{enumerate}

\textbf{(b)} Fünf Regeln (Beispiele): aussagekräftiger Betreff; passende Anrede und
Grußformel; kurz und klar formulieren; Empfänger bewusst wählen (CC/BCC sparsam);
Rechtschreibung prüfen und Anhänge ankündigen/beifügen; Signatur mit Pflichtangaben.

\textbf{(c)} Geschäftliche E-Mails ersetzen den postalischen Geschäftsbrief und gelten als
geschäftliche Korrespondenz. Deshalb müssen sie --- geregelt durch das \emph{EHUG} --- die
gleichen \emph{Pflichtangaben} enthalten wie ein Geschäftsbrief, damit der Empfänger im
gleichen Umfang über den Absender informiert wird. Der genaue Umfang hängt von der
\emph{Rechtsform} des Unternehmens ab.

\textbf{(d)} Ein Messenger/Chat ist ungeeignet für verbindliche, rechtlich relevante oder
vertrauliche Inhalte, die dokumentiert und beweissicher sein müssen (z.\,B.\ Verträge,
Bestellungen, Kündigungen) --- dafür ist die schriftliche Korrespondenz (E-Mail/Brief)
vorzuziehen.
\end{loesung}

\subsection*{Lösung Aufgabe~5}
\begin{loesung}
\textbf{(a)} Eine \emph{Gruppe} ist eine bloße Ansammlung von Personen ohne zwingend
gemeinsames Ziel. Ein \emph{Team} verfolgt ein \emph{gemeinsames Ziel}, die Mitglieder
ergänzen sich mit unterschiedlichen Rollen/Fähigkeiten und fühlen sich \emph{gemeinsam}
verantwortlich.

\textbf{(b)} Vorteile (drei): mehr Ideen/Kreativität, Ergänzung von Stärken/Arbeitsteilung,
höhere Akzeptanz von Entscheidungen, Motivation, gegenseitige Vertretung.
Nachteile (drei): höherer Abstimmungs-/Zeitaufwand, Konfliktpotenzial, soziales Faulenzen,
Gruppendruck/Gruppendenken.

\textbf{(c)} \emph{Soziales Faulenzen} (Trittbrettfahren): Einzelne reduzieren ihre
Leistung, weil sie sich auf das Team verlassen. Gegenmaßnahmen: klare Rollen- und
Aufgabenverteilung, messbare Einzelbeiträge, überschaubare Teamgröße, transparente Ziele.
\end{loesung}

\subsection*{Lösung Aufgabe~6}
\begin{loesung}
\textbf{(a)} Teamphasen nach Tuckman:
\begin{enumerate}[label=\arabic*., itemsep=0pt]
  \item \emph{Forming} (Orientierung): Team findet sich, höflich-unsicher, Ziele/Rollen
    unklar.
  \item \emph{Storming} (Konflikt): Auseinandersetzungen über Rollen, Aufgaben und Einfluss.
  \item \emph{Norming} (Organisation): gemeinsame Regeln und Rollen, Zusammenhalt wächst.
  \item \emph{Performing} (Leistung): effizientes, selbstorganisiertes Arbeiten am Ziel.
  \item \emph{Adjourning} (Auflösung): Abschluss, Auswertung, Auflösung.
\end{enumerate}

\textbf{(b)} Das Team befindet sich in der \textbf{Storming-Phase}: Es kommt zu Streit
über die Aufgabenverteilung, und einzelne Mitglieder ziehen sich zurück --- das ist
typisch für die Konfliktphase.

\textbf{(c)} Feedbackregeln (vier): zeitnah und konkret auf beobachtbares Verhalten
bezogen; Ich-Botschaften statt Du-Vorwürfe; beschreibend statt bewertend; wertschätzend
und konstruktiv mit Veränderungsvorschlag; der Empfänger lässt ausreden und rechtfertigt
sich nicht sofort.

\textbf{(d)} Die Aussage ist falsch. Die Storming-Phase mit Konflikten ist ein
\emph{normaler und notwendiger} Schritt der Teamentwicklung. Werden die Konflikte
konstruktiv gelöst, folgt die Norming- und Performing-Phase --- das Team wird gerade
\emph{dadurch} leistungsfähig.
\end{loesung}

\end{document}
